\chapter{Theoretical Background}\label{ch:theoretical_background}

The gap identified in \cref{ch:related_work}, that no existing tool composes style rewrites with data-level pruning into a joint pipeline, sits at the intersection of four traditions that each supply analogues for one layer of our pipeline.
We treat map rendering architecture (\cref{sec:map_rendering_approaches}), compiler optimization (\cref{sec:compiler_optimizations}), query optimization (\cref{sec:query_optimizations}), and columnar encoding (\cref{sec:data_encoding}) in turn, then close with the benchmarking methodology that grounds the empirical chapter (\cref{sec:benchmarking}).

\section{Map Rendering Approaches}\label{sec:map_rendering_approaches}

When building a slippy map (i.e.\ a map with the modern pan, zoom, rotation, and interaction capabilities), the technology stack is fairly simple.

A user interacts with an input device (mouse, touch, keyboard) to control the map view.
The map view is rendered using one of the following approaches:

\begin{description}
\item[\textbf{\ac{SVG}}] \begin{sloppypar} is a vector graphics format that can be used to render maps as shown by the ID Editor\footnote{\url{https://github.com/openstreetmap/iD}, last accessed 2026-03-16} or Leaflet\footnote{\url{https://leafletjs.com/}, last accessed 2026-03-16}.
Using \acp{SVG} means one is limited to a 2D rendering approach due to not having a camera.\end{sloppypar}

\item[\textbf{Canvas 2D}] is the HTML5 \texttt{<canvas>} immediate-mode drawing \ac{API}.
OpenLayers\footnote{\url{https://openlayers.org/}, last accessed 2026-04-21} uses Canvas~2D as its primary renderer, rasterising vector features on the \ac{CPU} rather than submitting them to the \ac{GPU}.
This avoids the complexity of \ac{GPU} pipeline management and shader compilation, but rendering throughput is \ac{CPU}-bound and does not benefit from hardware-accelerated geometry processing.
Like \ac{SVG}, Canvas~2D is limited to 2D views without pitch or rotation.

\item[\textbf{Low-level 3D Graphics \acp{API}}] \begin{sloppypar} like \ac{WebGL}\footnote{\url{https://get.webgl.org/}, last accessed 2026-03-16}, \ac{WebGPU}~\cite{WebGPU}, OpenGL\footnote{\url{https://registry.khronos.org/OpenGL-Refpages/gl4/}, last accessed 2026-03-16}, Metal\footnote{\url{https://developer.apple.com/metal/}, last accessed 2026-03-16}, Vulkan\footnote{\url{https://www.vulkan.org/}, last accessed 2026-03-16}, DirectX\footnote{\url{https://microsoft.github.io/DirectX-Specs/}, last accessed 2026-03-16}, etc.\ for rendering 3D graphics as shown by mapbox\footnote{\url{https://www.mapbox.com/}, last accessed 2026-03-16}, maplibre\footnote{\url{https://maplibre.org/}, last accessed 2026-03-16}, deck.gl\footnote{\url{https://deck.gl/}, last accessed 2026-03-16}, here\footnote{\url{https://www.here.com}, last accessed 2026-03-16} and many more.
This allows a more complex 3D-based rendering with pitch, yaw, and roll.
Due to having better hardware access, these \acp{API} can render more complex scenes with better performance.\end{sloppypar}
\end{description}

Due to the more advanced features and mostly better performance, we focus exclusively on renderers based on open-source 3D graphics \acp{API}.
While these renderers are powerful, they require data to be preprocessed and optimized specifically for rendering.
As a result, they cannot directly render high-level primitives such as shapes (polygons, lines, points), text, or images.
Instead, they operate on lower-level graphics primitives such as textures, triangles, and quads.

These low-level graphics building blocks are linked to map data through sprites, fonts, and styles, which are then interpreted by a rendering engine.

\begin{figure}[htbp]
	\centering
	\begin{tikzpicture}[
		node distance=0.8cm and 1cm,
		box/.style={rectangle, draw, rounded corners=2pt, minimum height=0.7cm,
			minimum width=2.0cm, font=\small, align=center},
		renderer/.style={box, fill=roleGen!20, minimum width=4.5cm, minimum height=0.7cm},
		server/.style={box, fill=roleSrc!20, minimum width=4.5cm, minimum height=0.7cm},
		datasource/.style={box, fill=roleData!20, minimum width=4.5cm, minimum height=0.7cm},
		arrow/.style={-Stealth, thick},
		label/.style={font=\scriptsize\itshape, fill=white},
		]
		% Map Renderer (top)
		\node[renderer] (renderer) {Map Renderer};
		
		% Tile Server
		\node[server, below=1cm of renderer] (server) {Tile Server};
		
		% Arrows from server to renderer with resource labels
		\coordinate (s1) at ($(server.north)!0.25!(server.north east)$);
		\coordinate (s2) at ($(server.north)!0.75!(server.north east)$);
		\coordinate (s3) at ($(server.north)!0.25!(server.north west)$);
		\coordinate (s4) at ($(server.north)!0.75!(server.north west)$);
		
		\draw[arrow] (s4) -- node[label] {style}   (s4 |- renderer.south);
		\draw[arrow] (s3) -- node[label] {tiles}   (s3 |- renderer.south);
		\draw[arrow] (s1) -- node[label] {sprites} (s1 |- renderer.south);
		\draw[arrow] (s2) -- node[label] {fonts}   (s2 |- renderer.south);
		
		% "over the network" brace
		\draw[decorate, decoration={brace, amplitude=6pt}, thick]
		([xshift=4pt]renderer.south east) -- ([xshift=4pt]server.north east)
		node[midway, right=10pt, font=\scriptsize\itshape, align=left]
		{over the network};
		
		% Data Source (bottom)
		\node[datasource, below=1cm of server] (datasource) {Data Source};
		
		\draw[arrow] (datasource.north) -- node[label] {convert to serving-friendly representation} (server.south);
	\end{tikzpicture}
	\caption{
Basic building blocks of a map rendering stack.
Data is taken from a data source, converted into a simplified format (for example a spatial index), and served to a client.
The visual appearance is defined by a style and supported by fonts and sprites for text and icons.}
	\label{fig:basic-arch}
\end{figure}

As a lower-level representation layer for displaying map content, the following components are required, as shown in \cref{fig:basic-arch}.
\emph{Styles} define the visual appearance of the map and reference the other required resources.
\emph{Tiles} provide the geometry (points, lines, polygons) and associated attributes (for example, building heights) displayed on the map.
\emph{Sprites} supply the icons.
\emph{Fonts} provide the glyphs used to render text labels.
Occasionally, additional assets such as 3D geometry models or terrain are used as well.

There are several approaches to deliver client-side rendered vector maps.

Styles, sprites, and fonts are typically served from a file server or \ac{CDN} and require little processing.
In some cases, parts of these resources may be generated dynamically, for example to support complex text shaping for Arabic, Hebrew, Indic, or Japanese scripts as described by Wipfli~\cite{wipfliComplexTextPlugin2025} for MapLibre or Lengyel~\cite{lengelDecadeSlug} for SLUG.
This processing is currently usually handled ahead of time on the server-side.

Tiles, by contrast, can be stored and accessed using three major approaches:

\begin{description}
\item[\textbf{Dynamic}]
    Store geospatial data in a database (for example, PostGIS\footnote{\url{https://postgis.net/}, last accessed 2026-05-23}) and use a tile server to generate and serve tiles on demand.
\item[\textbf{Client-driven}]
    Store encoded tiles as a single file in a blob store using a format that includes an internal tile index (for example, PMTiles or VersaTiles).
    Clients retrieve tiles using multiple \ac{HTTP} requests with appropriate Range headers to locate the required data.
\item[\textbf{Semi-static}]
    Store encoded tiles in a database file on disk (for example, SQLite) and use a tile server to extract and serve tiles to clients.
\end{description}

Dynamic data sources can be updated frequently by applying incremental diffs, whereas the other approaches typically require a full rebuild when the underlying data changes.
At larger scales, in-memory caching of tiles via regional or \ac{CDN}-style infrastructure, combined with fallback to runtime-dynamic data, is essential for performance and cost.

Although all of these resources can be optimized, tiles are usualy the most critical due to their continuously streaming nature as a user moves the map.
They usualy have a cumulatively higher data volume due to this, even though their individual tiles may be much smaller.
All other resources may depend on locale slightly (for example loading different Unicode ranges or fonts), but are only handled one time.

Real-world observations support this:
the performance of highly optimized systems such as Microsoft Bing Maps differs significantly from less optimized deployments, such as basemap.de.

In the following sections we examine each resource in detail, focusing on the performance characteristics and optimization strategies relevant to our work:

\subsection{Fonts}

Fonts can be client rendered from a font file.
This requires shipping a text shaping engine to the client, as neither native nor web targets have for example Harfbuzz\footnote{\url{https://github.com/harfbuzz/harfbuzz}, last accessed 2026-05-26} exposed as a feature.
This also means that for some fonts up to \qty{131}{\mega\byte}\cite{FontThatIncludes} of font data needs to be transferred.

For interactive use cases, this is prohibitive as the time to download this would be longer than acceptable.
Some fonts need this shaping as they need to connect neigboring characters neatly, but Latin fonts do usually not.
For latin-like characters, there is usually a clear 1:1 relationship between the Unicode code point and a rendered character.

This is why a simpler way exists:
Cutting fonts up into Unicode slices (\texttt{start..end}) and serving only the ranges that are needed to the client.
This cutting of fonts has the implicit advantage of pruning the space of possible font data to just the data that is needed to render the map.

Common open-source renderers like MapLibre or source-available renderers like Mapbox do this by shipping pre-rasterised glyph ranges as \acp{PBF} encoded \acp{SDF}~\cite{greenImprovedAlphatestedMagnification2007}, indexed by Unicode range (e.g.\ \texttt{0-255.pbf}, \texttt{256-511.pbf}).

\subsection{Sprites}

Sprite rendering is a well-established technique in computer graphics.
To avoid having to request hundreds to thousands of small files or to have to render SVGs on a client, this work is shifted to the server and the server delivers sprite sheets.
This allows efficient \ac{GPU} Sprite instancing.

Common renderers like MapLibre or Mapbox do this by shipping plain \ac{PNG} or \acs{SDF}-\acs{PNG}~\cite{greenImprovedAlphatestedMagnification2007} sprite sheets together with a lookup table.
The lookup table records the location of each sprite, whether it is an SDF, and optional stretchable regions for symbols such as highway shields.

\subsection{Tiles}\label{sec:tiles}

Tiles are the primary data delivery mechanism for interactive web maps.
Rather than transferring the entire dataset to the client, the map is partitioned into a regular grid of square tiles at discrete zoom levels following the XYZ tiling scheme.

At zoom level~$z$, the world is subdivided into $2^z \times 2^z$ tiles, each identified by its column~$x$, row~$y$, and zoom level~$z$.
Lower zoom levels cover large geographic areas with simplified geometry with higher zoom levels cover smaller areas with greater detail.
The client requests only the tiles intersecting the current viewport, and as the user pans or zooms, new tiles are fetched while old ones are discarded or cached.
When a client requests a tile at a zoom level beyond the source's maximum, the renderer \emph{overzooms}: it upscales the nearest available tile from the source's highest available zoom to fill the viewport.

The \emph{layers} sourced from that tile remain rendered at the overzoomed level even though no new tile is fetched.
A layer's effective visibility at a given camera zoom is therefore a function of both the layer's filter and the source's \mlprop{maxzoom}, not the layer bounds alone.
This observation has direct consequences for any optimizer that rewrites zoom bounds.
Removing a layer at zoom~14, for example, will suppress its rendering at zoom~15 and beyond if the source's \mlprop{maxzoom} is~14, because the renderer continues to display overzoomed tiles from the highest available source zoom.
We turn this into the asymmetric \mlprop{minzoom}/ \mlprop{maxzoom} tightening rule in \cref{sub:structural_optimizations}.
\Cref{lst:overzoom} shows the two operations the slicer performs to synthezise a target tile from the source's \mlprop{maxzoom} tile.
\stepcirc{1}~Rescale each coordinate by $2^{\Delta z}$ relative to the target-tile offset ($offsetX \notin \{null,0\} \lor offsetY \notin \{null,0\} \Rightarrow x = x * scale - offsetX$). \stepcirc{2}~Clip to the target extent with a 128-unit overdraw buffer for continuous edge geometry.
Both operations consume the source's \mlprop{maxzoom} tile as input, which is what makes \mlprop{maxzoom} tightening risky in a way that \mlprop{minzoom} tightening is not.

\clearpage
\begin{lstlisting}[
  language=TypeScript,
  caption={Condensed MapLibre GL-JS overzoom implementation (\texttt{src/source/vector\_tile\_overzoomed.ts})},
  label=lst:overzoom,
  escapeinside={(*@}{@*)},
]
function sliceVectorTileLayer(sourceLayer, maxZoomTileID, targetTileID) {
    const {extent} = sourceLayer;
    const dz = targetTileID.z - maxZoomTileID.z;
    const scale = Math.pow(2, dz);

    // Target tile's offset within the source tile, in target coords
    const offsetX = (targetTileID.x - maxZoomTileID.x * scale) * extent;
    const offsetY = (targetTileID.y - maxZoomTileID.y * scale) * extent;

    for (const feature of sourceLayer) {
        let geometry = feature.loadGeometry();
        // (*@\stepcirc{1}@*) Rescale each coordinate into the target tile's space
        for (const ring of geometry) for (const p of ring) {
            p.x = p.x * scale - offsetX;
            p.y = p.y * scale - offsetY;
        }
        // (*@\stepcirc{2}@*) Clip to the target extent with a 128-unit overdraw buffer
        geometry = clipGeometry(geometry, feature.type,
                                -128, -128, extent + 128, extent + 128);
        // ...emit feature with rescaled+clipped geometry
    }
}
\end{lstlisting}

No symmetric \enquote{underzoom} behavior exists:
If a source has no tile at a zoom lower than \mlprop{minzoom}, the layer is  not rendered there.

A \emph{vector tile} encodes map data in a compact binary format rather than in a pre-rendered image.
Each tile is organized into \emph{source layers} (for example, roads, buildings, water), and each source layer contains a set of \emph{features}.
A feature consists of geometry (point, linestring, or polygon encoded as integer coordinates relative to the tile's extent), optionally a set of key-value \emph{properties} (for example, \texttt{class=motorway} or \texttt{name=Main~Street}), optionally a set of IDs, and metadata such as the tile extent.

The following formats are used for encoding and transferring data to the client:

\begin{description}
    \item[\textbf{Raster tiles}] deliver pre-rendered images (\ac{PNG}, JPEG, WebP, Avif, JPEGXL) rather than vector geometry.
    They require no client-side rendering logic but cannot be restyled, rotated at arbitrary angles (beyond the affine transform imposed by globe mode), or queried interactively, as the visual representation is fixed at generation time.

    \item[\textbf{Analysis-focused formats}] such as FlatGeobuf\footnote{\url{http://flatgeobuf.org/}, last accessed 2026-05-11}, GeoParquet\footnote{\url{https://geoparquet.org/}, last accessed 2026-05-11}, and GeoArrow are optimized for analytical queries (filtering, aggregation, joins) rather than tile-based rendering.
    GeoParquet's row-group predicate pushdown and column pruning are analogous to the tile shaving we apply to \ac{MLT}, but they operate at per-file granularity.
    Parquet's row-group and column-chunk metadata would dominate at typical tile sizes (\qtyrange{1}{10}{\kilo\byte}), so we do not consider these formats further.

    \item[\textbf{\ac{GeoJSON}}]~\cite{butlerGeoJSONFormat2016} is a \ac{JSON}-based geospatial interchange format.
    Its human-readable structure makes it useful for debugging and small overlays, but its verbosity (text-encoded coordinates and repeated property keys per feature) makes it impractical for the continuous tile streaming required by interactive map renderers.

    \item[\textbf{\acf{MVT}}]\footnote{\url{https://github.com/mapbox/vector-tile-spec/tree/master/2.1}, last accessed 2026-04-21} is the dominant vector tile format, adopted as a \emph{de~facto} standard across open-source map renderers including MapLibre.
    It uses the \ac{PBF}\footnote{\url{https://protobuf.dev/}, last accessed 2026-04-21} as its wire format and stores features in a \emph{row-oriented} layout:
\begin{itemize}
    \item Each \texttt{Feature} carries its own list of property tags encoded as alternating key and value indices into per-layer string and value tables.
    \item Geometry is encoded as a sequence of \texttt{MoveTo}, \texttt{LineTo}, and \texttt{ClosePath} commands with delta-encoded integer coordinates relative to the tile extent (typically \texttt{4096} units per tile edge).
    This is also known as "turtle style graphics"\footnote{\url{https://docs.python.org/3/library/turtle.html}, last accessed 2026-5-27} as an analogy to the \texttt{turtle} python module.
    \item The row-oriented layout means that accessing a single property column requires parsing past all other properties of the same \texttt{Feature}, and that compression operates on the interleaved byte stream rather than on homogeneous data columns.
\end{itemize}


    \item[\textbf{\acf{MLT}}] \begin{sloppypar} is a columnar vector tile format proposed by Tremmel et~al.~\cite{tremmelMapLibreTileNext2025b}.
    Instead of the row-major key-value-pair layout of \ac{MVT}, \ac{MLT} stores each property as an independent column with type-specific encoding.
    This columnar layout enables per-column compression strategies (for example, dictionary encoding for low-cardinality string columns, delta encoding for monotonically increasing integers, or run-length encoding for repeated values) and eliminates the per-feature tag overhead of the protobuf encoding.
    Tremmel et~al.\ report up to $3\times$ tile size reduction on encoded tilesets and up to $6\times$ on certain large tiles, with the realised ratio depending on data characteristics.\end{sloppypar}
\end{description}

Beyond the choice of tile format, a number of transfer-layer optimizations can improve delivery performance:
utilising the approximately \qty{14}{\kilo\byte} TCP slow-start window as discussed by Nathaniel~\cite{WhyYourWebsite}, edge caching\footnote{such as suggested by \url{https://docs.here.com/map-rendering/docs/cache-maps-raster-tile-api}, last accessed 2026-03-16}, HTTP/2 multiplexing~\cite{belsheHypertextTransferProtocol2015} (and its QUIC-based successor HTTP/3~\cite{bishopHTTP32022} for lossy mobile networks), and modern content-encoding algorithms such as Zstandard~\cite{colletZstandardCompressionApplication2021} and Brotli~\cite{alakuijalaBrotliCompressedData2016} that achieve higher compression ratios than gzip's LZ77 foundation~\cite{zivUniversalAlgorithmSequential1977}.
Tile archive formats such as PMTiles~\cite{PMTilesConceptsProtomaps} and COMTiles~\cite{tremmelCOMTILESCASESTUDY2023} address the container layer by enabling serverless, cloud-native tile delivery from a single file in blob storage.
Since \ac{MVT} is the established baseline and \ac{MLT} aims to improve upon it, we evaluate both formats and use the transition from \ac{MVT} to \ac{MLT} as one of our data-level optimization strategies.

\subsection{Styles}\label{sec:styles}

A map style defines the complete visual appearance of a map: which data is shown, how it is colored, sized, and labelled, and at which zoom levels each element appears.
Three architectural approaches to map styling have been established:

\begin{description}
    \item[\textbf{Server-side declarative}] styles, such as MapCSS, define rendering rules on the server.
    The server evaluates the style against the data and produces pre-rendered raster tiles that the client displays without further interpretation.
    This simplifies the clients rendering logic, but prevents client-side restyling and thus limits interactivity.

    \item[\textbf{Client-side declarative}] \begin{sloppypar} styles, as used by MapLibre\footnote{\url{https://maplibre.org/}, last accessed 2026-03-16} and Mapbox\footnote{\url{https://www.mapbox.com/}, last accessed 2026-03-16}, transmit a machine-readable style document to the client, which interprets the style at render time.
    This style can also be shipped with the code of the rendering engine to the client or partially constructed (e.g. road shields) there.
    This enables dynamic restyling, data-driven visualization, and interactive features such as hover effects and click queries without regenerating tiles.\end{sloppypar}

    \item[\textbf{Callback-based}] styling, as proposed by Gleo\footnote{\url{https://ivansanchez.gitlab.io/gleo/}, last accessed 2026-03-16} or Navara\footnote{\url{https://navara-docs.netlify.app/}, last accessed 2026-04-21}, replace declarative rules with imperative callbacks that the application registers with the renderer.
    This offers maximum flexibility but makes static analysis and automatic optimizations impossible, since the styling logic is opaque to the rendering engine.
\end{description}

We focus on client-side declarative styles, specifically the MapLibre style specification.
It is the most widely adopted approach in open-source map rendering, and its declarative nature makes it amenable to static analysis and automatic optimizations.

The MapLibre style specification\footnote{\url{https://www.maplibre.org/maplibre-style-spec/}, last accessed 2025-10-24} defines a map's appearance as a \ac{JSON} document with the following key components:

\begin{description}
    \item[\textbf{Sources}] declare where tile data comes from (tile server URL, or inline \ac{GeoJSON}) and at which zoom levels tiles are available (\mlprop{minzoom}, \mlprop{maxzoom}).

    \item[\textbf{Layers}] are the central rendering unit.
    Each layer references a source and a \emph{source-layer} within that source, specifies a \emph{type} (\texttt{\mlval{fill}}, \texttt{\mlval{line}}, \texttt{\mlval{circle}}, \texttt{\mlval{symbol}}, \texttt{\mlval{fill-extrusion}}, \texttt{\mlval{raster}}, \texttt{\mlval{background}} or \texttt{\mlval{hillshade}}) and carries a set of paint and layout properties that control its appearance.
    Layers are drawn in declaration order from bottom to top.
    \begin{description}

    \item[\textbf{Filters}] are boolean expressions attached to layers that select which features from the source-layer are rendered.
    A feature is drawn by a layer only if the layer's filter evaluates \mlval{true} for that feature's properties and geometry type.

    \item[\textbf{Paint and layout properties}] control the visual appearance and placement of rendered features.
    Paint properties (colors, opacities, widths) can be changed without triggering a re-layout.
    Layout properties (text anchoring, icon rotation, symbol spacing) affect feature placement and require re-evaluation when modified.

    \item[\textbf{Metadata bounds}] (\mlprop{minzoom}, \mlprop{maxzoom}) restrict the zoom range in which a layer is active.
    This allows the renderer to skip layers outside that range before filter evaluation or geometry decoding.
    \Cref{lst:ishidden} shows the corresponding \texttt{isHidden} guard.
    The worker invokes \texttt{layer.isHidden(this.zoom, true)} to skip bucket creation during tile parsing, while the renderer returns early from \texttt{renderLayer} on \texttt{layer.isHidden(this.transform.zoom)} before dispatching draw calls.
    
    \begin{minipage}{\dimexpr\linewidth\relax}
\begin{lstlisting}[language=TypeScript,caption={MapLibre GL-JS \texttt{isHidden} method from \texttt{src/style/style\_layer.ts}},label=lst:ishidden]
isHidden(zoom: number = this.minzoom, roundMinZoom: boolean = false) {
    if (this.minzoom && zoom < (roundMinZoom ? Math.floor(this.minzoom) : this.minzoom)) return true;
    if (this.maxzoom && zoom >= this.maxzoom) return true;
    return this._evaluatedVisibility === 'none';
}
\end{lstlisting}
\end{minipage}

\Cref{lst:ishidden_usage} reproduces the two call-sites that establish those two stages: the worker calls \texttt{layer.isHidden(this.zoom, true)} and \texttt{continue}s before bucket creation, and the painter's \texttt{renderLayer} returns immediately on \texttt{layer.isHidden(this.transform.zoom)} before dispatching any \texttt{draw.fill}/ \texttt{draw.line}/ \texttt{draw.symbol} call.
\end{description}
\end{description}

Filters and properties can be specified as literal constants, as \emph{zoom-dependent} expressions (varying with the camera zoom level via \mlop{interpolate} or \mlop{step} ramps), as \emph{data-driven} expressions (varying per feature based on property values via \mlop{match}, \mlop{case}, or \mlop{get} operators), or as a combination of both.
The resulting expression trees can be deeply nested and may contain redundant or unreachable branches.
These characteristics make them amenable to the same kinds of simplification and optimizations techniques used in compilers and database query planners.

\Cref{lst:worker_pipeline} shows the worker's tile-parse pipeline as five named steps, marked inline as \stepcirc{1} to \stepcirc{5}.

\noindent Every optimization we present in \cref{ch:methods} attaches to one of these steps.
Merging layers reduces work for \stepcirc{1}.
Pruning features server-side makes \stepcirc{2} not handle them.
Tightening \mlprop{minzoom}/ \mlprop{maxzoom} makes the \stepcirc{3} guard true earlier.
Pre-folding zoom-dependent constants reduces the work in \stepcirc{4}.
Filter-to-property propagation simplifies the per-feature work inside \stepcirc{5}.

\begin{lstlisting}[
    language=TypeScript,
    escapeinside={(*@}{@*)},
    caption={MapLibre GL-JS WorkerTile.parse processing pipeline (\texttt{src/source/worker\_tile.ts})},
    label=lst:worker_pipeline
]
// (*@\stepcirc{1}@*) Iterate source layers
const layerFamilies = layerIndex.familiesBySource[this.source];
for (const sourceLayerId in layerFamilies) {
    const sourceLayer = data.layers[sourceLayerId];
	if (!sourceLayer) continue;

    ...

    const sourceLayerIndex = sourceLayerCoder.encode(sourceLayerId);

    // (*@\stepcirc{2}@*) Collect raw features
    const features = [];
    for (let index = 0; index < sourceLayer.length; index++) {
        const feature = sourceLayer.feature(index);
        const id = featureIndex.getId(feature, sourceLayerId);
        features.push({feature, id, index, sourceLayerIndex});
    }

    for (const family of layerFamilies[sourceLayerId]) {
        const layer = family[0];

        ...
        // (*@\stepcirc{3}@*) Skip hidden layers
        if (layer.isHidden(this.zoom, true)) continue;

        // (*@\stepcirc{4}@*) Recalculate zoom-dependent expressions
        recalculateLayers(family, this.zoom, availableImages);

        const bucket = buckets[layer.id] = layer.createBucket({...});

        // (*@\stepcirc{5}@*) Filter features and build vertex data
        bucket.populate(features, options, this.tileID.canonical);
        featureIndex.bucketLayerIDs.push(family.map((l) => l.id));
    }
}
\end{lstlisting}

\subsection{Feature and Global State}\label{sub:feature_state}

A MapLibre map carries two runtime key-value stores that sit alongside the JSON style document and the encoded tile data, even though they live in neither.
\emph{Feature state} is a per-feature map keyed by the (source, source-layer, feature ID) triple, mutated from the host application via \texttt{setFeatureState} and \texttt{removeFeatureState}, and read from style expressions through the \texttt{[\mlop{"feature-state"},~key]} operator.
Its canonical use is interactive UI such as hover highlighting or click-driven selection.
\emph{Global state} is a map-wide key-value store mutated via \texttt{setGlobalStateProperty} and read through the \texttt{[\mlop{"global-state"},~key]} operator.
It typically backs up application-wide toggles such as a day/night switch or a per-mode road-class filter.

A state change re-evaluates the dependent paint expressions at render time, without re-parsing the tile or re-running bucket population.
This is what keeps feature state cheap enough to drive hover effects.
Both stores are mutated by the host application at runtime, after the style document and tile data have been loaded, so their values are not determined by either.
A reference whose key is missing from the store, or whose stored value has a type the consuming expression cannot handle, raises a runtime evaluation error that aborts rendering for the affected expression.
Feature-state addressing further requires stable feature IDs, since the keying triple includes the feature ID.

\section{Compiler Optimization Techniques}\label{sec:compiler_optimizations}

The style expressions used by map renderers such as MapLibre are small programs in a special purpose \ac{DSL}:
they contain conditionals (\mlop{case}, \mlop{match}), boolean operators (\mlop{all}, \mlop{any}), arithmetic, and data access (\mlop{get}, \mlop{has}), and variables (lexical bindings via \mlop{let}/ \mlop{var}, and the runtime stores of \cref{sub:feature_state}).
Optimising these expressions draws directly on techniques developed over decades in compiler construction~\cite{ahoCompilersPrinciplesTechniques2007}.
In this section we introduce the foundational concepts we apply in \cref{ch:methods}.

\subsection{Peephole Optimization}

\begin{sloppypar}
A \emph{peephole optimizer} examines a small, fixed-size window of code (the \enquote{peephole}) and applies local pattern-matching rewrites that replace instruction sequences with shorter or faster equivalents~\cite{davidsonCodeSelectionObject1984}.
Each rewrite rule is independent, matching a local pattern and emitting a replacement without requiring knowledge of the surrounding program.
Many small, local improvements compose into significant global optimizations under repeated application.
\end{sloppypar}

Peephole optimizers are classified by the level at which they operate.
\emph{Machine-level} peephole optimizers replace instruction sequences (e.g., replacing a multiply by a power of two with a shift).
\emph{\ac{IR}-level} peephole optimizers apply algebraic simplifications, identity eliminations, and strength reductions to expression trees.
In the context of style expressions, the peephole window corresponds to a single expression node and its immediate children.
A post-order visitor walks the expression tree and attempts to apply each rule at every node, producing a locally simplified tree.

\subsection{Constant Folding and Propagation}

\emph{Constant folding} evaluates expressions whose operands are all compile-time constants, replacing the expression with its result~\cite{ahoCompilersPrinciplesTechniques2007}.
For example, the expression $3 + 5$ is replaced by $8$.
In a style context, a comparison such as \texttt{[\mlop{"=="},~\mlval{2},~\mlval{3}]} can be evaluated at optimizations time and replaced by \mlval{false}.
\emph{Constant propagation} extends this idea by tracking the known values of variables through the program's control flow and substituting them at their use sites.
If a variable is assigned a constant, all subsequent references to that variable can be replaced by the constant, potentially enabling further constant folding at the use sites.

A more powerful variant, \acf{SCCP}, combines constant propagation with dead branch analysis in a single pass.
It simultaneously determines which values are constant \emph{and} which branches are reachable, enabling the discovery of constants that are only visible when unreachable code is ignored~\cite{wegmanConstantPropagationConditional1991}.

\subsection{Dead Code Elimination}

\acf{DCE} removes program statements whose results are never used or whose control-flow paths are unreachable~\cite{ahoCompilersPrinciplesTechniques2007}.
Two common forms are \emph{unreachable code elimination}, which removes code that can never be executed (e.g., code after a guaranteed return, or branches guarded by a statically false condition), and \emph{dead store elimination}, which removes assignments to variables that are never subsequently read.

In the map-style domain, dead code elimination corresponds to removing style layers that can never produce visible output, because their filter is unsatisfiable, their geometry type does not match the data, or their paint properties render them invisible (e.g., zero opacity at all zoom levels).

\subsection{Algebraic Simplification}

\emph{Algebraic simplification} applies mathematical identities to reduce expression complexity without changing semantics~\cite{ahoCompilersPrinciplesTechniques2007}.
Implemented algebraic rules include:

\begin{itemize}
	\item \emph{Identity elements}:
	\[
	\begin{alignedat}{4}
		x \land \mathit{true}   & = x \qquad &
		x \lor \mathit{false}   & = x \qquad &
		x + 0                   & = x \qquad &
		x \cdot 1               & = x
	\end{alignedat}
	\]

	\item \emph{Annihilation}:
	\[
	\begin{alignedat}{3}
		x \land \mathit{false}  & = \mathit{false} \qquad &
		x \lor \mathit{true}    & = \mathit{true} \qquad &
		\forall x \in \mathbb{N}: x \cdot 0               & = 0
	\end{alignedat}
	\]

	\item \emph{De~Morgan's laws}:
	\[
	\begin{alignedat}{2}
		\lnot (a \land b) & = \lnot a \lor \lnot b \qquad &
		\lnot (a \lor b)  & = \lnot a \land \lnot b
	\end{alignedat}
	\]

	\item \emph{Idempotence}:
	\[
	\begin{alignedat}{2}
		x \land x & = x \qquad &
		x \lor x  & = x
	\end{alignedat}
	\]

	\item \emph{Absorption}:
	\[
	\begin{alignedat}{2}
		x \land (x \lor y) & = x \qquad &
		x \lor (x \land y) & = x
	\end{alignedat}
	\]
\end{itemize}

\noindent For style expressions, these rules simplify boolean filter expressions.
A filter \texttt{[\mlop{"all"}, [\mlop{"all"}, a, b], c]} is flattened to \texttt{[\mlop{"all"}, a, b, c]} by associativity.
A filter \texttt{[\mlop{"all"}, \mlval{true}, x]} is simplified to~\texttt{x} by identity elimination.

\subsection{Fixpoint Iteration}\label{sub:fixpoint}

Individual peephole rules may expose further opportunities.
Constant folding a subexpression may create a new constant parent that enables dead branch elimination, so a single traversal is generally insufficient.

The remedy is to apply the entire rule set repeatedly until no rule fires.
A term to which no rule applies is a \emph{fixpoint} of the rewrite function, equivalently a \emph{normal form} of the underlying rewrite system.
We treat our expression-level optimizer formally as a \ac{TRS}~\cite{baaderTermRewritingAll1999}, which makes the relevant theoretical lens the pair of \emph{termination}\footnote{every reduction sequence reaches a normal point} and \emph{confluence}\footnote{all sequnences reach the same one} rather than the lattice-theoretic monotone-convergence framing of classical dataflow fixpoint iteration~\cite{ahoCompilersPrinciplesTechniques2007}.
\Cref{sub:expression_optimizations} gives a lexicographic measure on which most rules strictly decrease, together with the two rules (De~Morgan and distributive factoring) that can temporarily increase it, so a Knaster-Tarski-style convergence theorem~\cite{tarskiLatticetheoreticalFixpointTheorem1955} does not apply.

In practice the loop reaches a normal form within a handful of iterations and we cap it at eight as a safety bound.
Newcomb et~al.~\cite{newcombVerifyingImprovingHalides2020} reported similar confluence and termination concerns for Halide's hand-written rewrite optimizer and used synthezis-based techniques to verify correctness and discover missing rules.

\section{Query Optimization}\label{sec:query_optimizations}

Style filters in map rendering are structurally equivalent to database query predicates:
both select a subset of records (features) from a dataset (tile) based on attribute conditions.
Techniques from relational query optimizations therefore apply directly.

\subsection{Cost-Based Query Optimization}

A \emph{cost-based query optimizer} selects among semantically equivalent query execution plans by estimating the cost of each plan and choosing the cheapest~\cite{selingerAccessPathSelection1979}.

The key insight is that the \emph{order} in which predicates are evaluated affects performance even when the final result is identical.
Under short-circuit semantics, a conjunction $p_1 \land p_2 \land \cdots \land p_n$ is fastest when the most-likely-\mlval{false} predicate comes first, and a disjunction $p_1 \lor p_2 \lor \cdots \lor p_n$ places the most-likely-\mlval{true} predicate first.

\subsection{Selectivity Estimation}

Estimating the expected fraction of records satisfying a predicate is known as \emph{selectivity estimation}.
Selinger et~al.~\cite{selingerAccessPathSelection1979} introduced the foundational formulas for estimating selectivity under the \emph{independence assumption}, the assumption that attribute values in different columns are statistically independent.

For a single equality predicate $A = v$ on a column with $d$ distinct values, the estimated selectivity is $s = 1/d$ (uniform distribution assumption).
When value-frequency histograms are available, the estimate improves to $s = \mathit{freq}(v) / N$, where $N$ is the total number of records.
Compound predicates are estimated by combining individual selectivities:
\begin{align*}\label{eq:selectivity_theory}
    s(p_1 \land p_2) &= s(p_1) \cdot s(p_2)\\
    s(p_1 \lor p_2) &= 1 - (1 - s(p_1))(1 - s(p_2))
\end{align*}

The independence assumption is known to be inaccurate when columns are correlated, but it provides a computationally cheap approximation that is effective in practice~\cite{selingerAccessPathSelection1979}.
More sophisticated estimators (such as multi-dimensional histograms, sampling-based methods, or the lightweight graphical models proposed by Tzoumas et~al.~\cite{tzoumasLightweightGraphicalModels2011} that capture inter-attribute correlations without the full cost of joint-distribution estimation) improve accuracy at higher computational cost.

Map-style filters are in fact known to violate independence:
In OpenMapTiles, for example, \texttt{class} and \texttt{subclass} are strongly correlated by schema design (every value of \texttt{subclass} implies a constrained set of \texttt{class} values), as are \texttt{admin\_level} and \texttt{maritime}.
A joint-distribution estimator, such as the one we propose in \cref{sec:selectivity_reordering}, would therefore give tighter bounds, but our selectivity signal is consumed only by a \emph{reordering} heuristic (not by a cost-based plan chooser) and reordering is monotone in the ordering of the individual selectivities rather than in their absolute values.
The independence-estimator error therefore propagates into the quality of the short-circuit ordering but never into the correctness of the rewritten style, which makes the cheap estimator an acceptable engineering choice for our use case.

\subsection{Partial Evaluation}\label{sub:partial_eval}

\emph{Partial evaluation} (also called program specialization) transforms a program with respect to known static inputs, producing a residual program that depends only on the remaining dynamic inputs~\cite{jonesPartialEvaluationAutomatic1993}.
The style optimizer performs partial evaluation of the rendering function: the style (static input) is known at optimization time, while tile data (dynamic input) varies at runtime.
Filter-to-property constant propagation (\cref{sub:structural_optimizations}) is a textbook instance.
MapLibre paints only features that pass the filter (\cref{lst:bucket_populate}), so any equality constraint in the filter holds for every feature reaching the paint expression.
The optimizer substitutes that constraint into the paint expression and folds away the branches it rules out, while the filter itself stays unchanged and still runs at render time against dynamic feature data.
The same framing covers the stats-driven folding pass, where the tile's value distribution (cardinalities, per-property frequencies, geometry-type counts) is the additional static input that lets the optimizer collapse \mlop{match} arms over absent values, tighten zoom bounds, and convert high-selectivity predicates into constants.
This is distinct from profile-guided optimization, which would require measuring per-expression hotness at render time and specialising the hot paths.
The style optimizer never observes render-time execution and therefore does not perform that kind of specialisation.

The encoding strategy competition (\cref{sub:encoding_strategies}) is no partial evaluation but a single-output search.
It trials each candidate from a designer-specified set (sort orders $\times$ encoding schemes) on the actual data and emits the smallest, with a profiler pruning the search beforehand.
We characterise this as \emph{exhaustive strategy selection} over a bounded candidate space, optimal within that space but blind to strategies outside it.

\section{Data Encoding and Compression}\label{sec:data_encoding}

Efficient encoding of tile data is essential for interactive map performance, as tiles are fetched continuously during map interaction.
A column has to be both \emph{represented} as a value sequence and \emph{packed} into bytes, and some named schemes only address the first step while relying on the second to actually shrink the byte count.

Following the MLT specification, we call the first kind a \emph{logical-level technique} and the second a \emph{physical-level technique}.
A column applies zero or more logical techniques in chain and then exactly one physical technique, for example DeltaRLE~$\rightarrow$~FastPFOR.

\Cref{sub:logical_integer_encodings,sub:physical_integer_compressions} catalogue the techniques we use in our data-level optimizations of \cref{ch:methods}.

\subsection{Column-Oriented Storage}\label{sub:column_storage}

Traditional database systems store data in \emph{row-oriented} layouts, where all attributes of a single record are stored contiguously in memory.
This layout is efficient for transactional workloads that access entire records, but wasteful for analytical queries that access only a few columns, as the entire record must be read to access each single attribute.

\emph{Column-oriented} (columnar) storage transposes this layout, storing all values of a single attribute contiguously~\cite{abadiDesignImplementationModern2013}.
The foundational system in this space is by Dremel~\cite{melnikDremelInteractiveAnalysis}, who introducsed the record shredding and assembly algorithm for nested columnar data using definition and repetition levels, the representation pattern that directly inspired Apache Parquet and, by extension, the columnar layout that \ac{MLT} adapts for tile data.
This layout offers three advantages for read-heavy workloads. \emph{Reduced I/O} follows from reading or transferring only the referenced columns. \emph{Better compression} follows from the higher local homogeneity of values of the same type and semantic domain compared to interleaved multi-type records. \emph{Vectorised processing} follows from the layout's natural fit for \ac{SIMD}-style batch operations over arrays of the same type.

In the context of vector tiles, the row-oriented \ac{MVT} format stores all properties of a feature together, whereas the columnar \ac{MLT} format stores each property as an independent column.
The columnar layout allows per-column encoding strategies:
a column of \texttt{road} classes can use dictionary encoding, while a column of building heights can use a DeltaRLE logical chain handed to a physical technique such as FastPFOR, each chosen independently.

\subsection{Logical Integer Encodings}\label{sub:logical_integer_encodings}

Integer columns appear throughout tile data: geometry coordinates, feature IDs, property values, and dictionary indices.
Logical-level techniques rewrite the value sequence so that a downstream physical technique (\cref{sub:physical_integer_compressions}) sees a distribution that packs to fewer bytes.
Several such techniques exploit structural properties of integer sequences:

\begin{description}
    \item[\textbf{Zigzag encoding}] maps signed integers to unsigned ones ($n \mapsto 2n$ for $n \geq 0$, $n \mapsto -2n-1$ for $n < 0$) so that small-magnitude negative values also produce small unsigned codes.
    The byte count is unchanged on its own and the saving comes from the downstream physical technique.

    \item[\textbf{Delta encoding}] replaces each value with the difference from its predecessor: given a sequence $v_0, v_1, \ldots, v_n$, the encoded sequence is $v_0,\, v_1 - v_0,\, v_2 - v_1,\, \ldots,\, v_n - v_{n-1}$.
    Like zigzag, delta does not shrink the column by itself.
    It produces the same number of values at (often) smaller magnitude, which then pack well under VarInt or FastPFOR.

    \item[\textbf{\acf{RLE}}] replaces consecutive runs of identical values with (value, count) pairs.
    Unlike zigzag and delta, RLE reduces the symbol count on its own.
    MLT still classifies it as logical because it composes with other logical steps before a physical technique.
    It is most effective for low-cardinality, clustered columns, for example a road-class column sorted by class.
\end{description}

Logical techniques chain freely.
The canonical chain applies delta first and then \ac{RLE} (written DeltaRLE), effective for approximately sorted sequences where delta produces long runs of small (often zero) values that \ac{RLE} then collapses.

\subsection{Physical Integer Compressions}\label{sub:physical_integer_compressions}

A physical-level technique packs a value sequence (possibly the output of a logical chain) into bytes, and each integer column applies exactly one.
Our encoder uses two:

\begin{description}
    \item[\textbf{\acf{VarInt}}] represents each integer with a variable number of bytes, using 7~bits per byte for data and 1~bit as a continuation flag, so values in $[0, 127]$ take 1~byte, $[128, 16383]$ take 2~bytes, and so on\footnote{\url{https://protobuf.dev/}, last accessed 2026-04-21}.
    VarInt expects unsigned values, which is why a signed integer column is normally preceded by a zigzag remap.

    \item[\textbf{\acf{FastPFOR}}] partitions the input into fixed-size blocks (typically 128 or 256 values), computes the minimum bit width that fits most values in each block, and packs at that width, with outliers in a separate patch list~\cite{lemireDecodingBillionsIntegers2015}.
    This achieves near-optimal bit packing when most values share a similar magnitude.
\end{description}

Choosing the physical technique for a given logical chain has been studied systematically.
Abadi et~al.~\cite{abadiIntegratingCompressionExecution2006} proposed property-based selection rules, and Damme et~al.~\cite{dammeComprehensiveExperimentalSurvey2019} developed a cost-model-driven strategy that picks the lowest expected cost for a given data distribution.
A third school sidesteps both rule and model and decides empirically: BtrBlocks~\cite{kuschewskiBtrBlocksEfficientColumnar2023} trial-encodes candidate schemes on a sample of each column block and keeps the smallest result, arguing that lightweight on-sample trials are both more accurate than a cost model and cheap enough to run at encode time.
Our encoding competition in \cref{sub:encoding_strategies} sits in the same trial-based family.
It trials a pruned set of (logical chain, physical technique) pairs and retains the smallest output, but operates on each whole tile rather than on a sample, since per-tile data volumes are modest.

\subsection{String Compression}\label{sub:string_compression}

String data in vector tiles includes feature names, road classes, land-use categories, and administrative identifiers.
The same logical-then-physical split applies as for integer columns.

\textbf{Dictionary encoding} is the logical-level string technique.
It stores each unique string once and replaces occurrences with integer indices, which is effective when many features share the same value (e.g. thousands of features with \texttt{class = \texttt{residential}}).
The resulting index column is itself handed to the integer techniques of \cref{sub:logical_integer_encodings,sub:physical_integer_compressions}.

\textbf{\acf{FSST}}~\cite{bonczFSSTFastRandom2020} is the physical-level string technique.
It is a lightweight, block-oriented compressor that learns a \emph{symbol table} of up to 255 frequent byte sequences from the input corpus and replaces occurrences with single-byte codes.
Unlike general-purpose compressors such as LZ77~\cite{zivUniversalAlgorithmSequential1977a} or ~\cite{colletZstandardCompressionApplication2021}, \ac{FSST} preserves random access, so each compressed string can be decompressed independently.
The symbol table is stored once and amortised over all strings in the column.

The two compose in the same logical-then-physical pattern.
Applying \ac{FSST} to the deduplicated dictionary strings compresses both the unique values and the byte-level redundancy within them.

\subsection{Space-Filling Curves}\label{sub:space_filling_curves}

\emph{Space-filling curves} map multi-dimensional points to a one-dimensional ordering while preserving spatial locality.
Points that are close in two-dimensional space tend to receive nearby positions on the curve~\cite{moonAnalysisClusteringProperties2001}.
Such curves play two roles in tile encoding.
As a logical encoding they map a 2D coordinate to a single integer column (\cref{sub:logical_integer_encodings}).
As a \emph{feature ordering} they sort features before per-column encoding so that spatially adjacent features end up contiguous, improving the effectiveness of delta and \ac{RLE} on geometry and property columns.
This subsection treats the feature-ordering use.

Two space-filling curves are commonly used.

The \textbf{Z-order (Morton) curve}~\cite{morton1966computer} interleaves the binary representations of the $x$ and $y$ coordinates to form a single index.
For coordinates $(x, y)$ with $b$-bit components, the Morton code is
\begin{equation}\label{eq:morton}
    M(x, y) = \sum_{i=0}^{b-1} \bigl( x_i \cdot 2^{2i} + y_i \cdot 2^{2i+1} \bigr),
\end{equation}
where $x_i$ and $y_i$ are the $i$-th bits of $x$ and $y$, respectively.
Morton codes are fast to compute via bit-interleaving instructions, but the Z-order curve exhibits \enquote{jumps} between quadrants that break spatial continuity.

The \textbf{Hilbert curve} traces a continuous path through the grid without the inter-quadrant jumps of the Z-order curve, achieving better locality preservation at the cost of a more complex index computation~\cite{moonAnalysisClusteringProperties2001}.
The Hilbert curve has also been used as the basis for spatial index structures: the Hilbert R-tree~\cite{kamelHilbertRtreeImproved1994} orders rectangles along the Hilbert curve to achieve tighter node packing than conventional R-trees, and FlatGeobuf\footnote{\url{http://flatgeobuf.org/}, last accessed 2026-05-11} uses a packed Hilbert R-tree for spatial indexing of its features.
The Hilbert index is computed recursively.
The unit square is divided into four quadrants, each mapped to a sub-range of the curve.
The quadrant containing the point determines the high-order bits of the index, and the process recurses on the sub-quadrant.

For tile data, we assign each feature a curve index based on its first vertex coordinate, and our encoding competition described in \cref{ch:methods} selects the feature ordering that produces the smallest encoded output (across Morton, Hilbert, and non-spatial orderings).

\subsection{Locality-Sensitive Hashing}\label{sub:lsh}

When multiple string columns share overlapping vocabularies, storing a single shared dictionary can reduce overhead compared to per-column dictionaries.
Identifying which columns share sufficient overlap is a set similarity problem.

\emph{MinHash}~\cite{broderResemblanceContainmentDocuments1998} is a locality-sensitive hashing technique that estimates the \emph{Jaccard similarity} between two sets $A$ and $B$:
\begin{equation}\label{eq:jaccard}
    J(A, B) = \frac{|A \cap B|}{|A \cup B|}
\end{equation}
MinHash yields no asymptotic improvement over the exact computation.
Constructing a signature requires examining every element of $A$ and $B$, so the time complexity remains $\mathcal{O}(|A|+|B|)$, and the input sets must be retained throughout construction.
The reduction is a constant factor and applies to the persistent representation.
Whereas the exact computation stores all $\mathcal{O}(|A|+|B|)$ elements, MinHash stores only $k$ hash values per set.
For fixed $k$, the saving in persistent storage grows without bound as $|A|$ and $|B|$ increase.
For each of $k$ independent hash functions $h_1, \ldots, h_k$, the minimum hash value over all elements is recorded: $\mathit{sig}_i(A) = \min_{a \in A} h_i(a)$.
The Jaccard similarity is then estimated as:
\begin{equation}\label{eq:minhash}
    \hat{J}(A, B) = \frac{1}{k} \sum_{i=1}^{k} \mathbf{1}\bigl[\mathit{sig}_i(A) = \mathit{sig}_i(B)\bigr]
\end{equation}
where $\mathbf{1}[\cdot]$ is the indicator function.
The estimate is unbiased with variance $\mathcal{O}(1/k)$~\cite{broderResemblanceContainmentDocuments1998}, so the accuracy is fine-tunable.


\section{Benchmarking}\label{sec:benchmarking}

Benchmarking refers to the systematic measurement of system performance using controlled workloads and standardized metrics.
The purpose of benchmarking is to evaluate system behavior, compare alternative implementations, and quantify the effects of optimizations.

According to Jain~\cite{jainArtComputerSystems1991}, a good evaluation must start with selecting relevant metrics.
Georges et~al.~\cite{georgesStatisticallyRigorousJava2007} later demonstrated that naive averaging of execution times on managed runtimes (such as JavaScript engines) leads to unreliable conclusions, advocating instead for confidence intervals computed over multiple independent \ac{JVM}/engine invocations.
That principle informs the bootstrap CI methodology we use.
The bootstrap itself, a resampling method for constructing confidence intervals without distributional assumptions, was formalised by Efron and Tibshirani~\cite{efronIntroductionBootstrap1998} and is the specific statistical technique we apply throughout our evaluation (\cref{ch:experiments}).
Hoefler and Belli~\cite{hoeflerScientificBenchmarkingParallel2015} codified twelve common pitfalls in reporting performance results.
Our evaluation in \cref{ch:experiments} follows their guidelines by reporting per-metric confidence intervals and separating throughput from latency measurements.
In the context of web maps, the candidate metrics span a range from deployment-critical on battery-constrained targets to directly user-visible ones on the desktop:

\begin{description}
    \item[\textbf{Resource utilisation,}] such as \ac{CPU} load, \ac{GPU} utilisation, and memory consumption to have pointers where further optimizations is possible.
    \item[\textbf{Latency and responsiveness,}] measuring the time between initiating an operation (for example zooming) and the first visible result.
    \item[\textbf{Throughput,}] typically measured in \acp{FPS} or frame time, describing how many frames can be rendered per second.
    \item[\textbf{Frame stability,}] which captures variations in frame time and identifies dropped frames or rendering stutters.
    \item[\textbf{Network utilisation,}] including transferred bytes, tile request rates, and bandwidth consumption.
    \item[\textbf{Power consumption}] in Watts, which is relevant for mobile devices and battery-powered systems.  We do not measure power consumption.
We use no \ac{RAPL} counters, battery deltas, or external power meters, as our benchmark host is a desktop workstation rather than a battery-constrained mobile device.  We list power here as a candidate metric because transfer-size reductions and reduced draw-call counts are established proxies for energy savings on mobile targets.
Cellular radio energy is proportional to data volume, and \ac{GPU} active time scales with draw-call count.
\end{description}

These can be benchmarked at different workload levels~\cite{jainArtComputerSystems1991}:

\begin{description}
    \item[\textbf{Micro-benchmarks}] isolate a single operation (e.g.\ decoding a single tile, evaluating a single filter expression) and measure its cost in isolation.
They are useful for identifying bottlenecks but may not reflect real-world performance due to caching and concurrency effects.
    \item[\textbf{Macro-benchmarks}] measure the performance of a complete subsystem (e.g.\ loading a full style and rendering a fixed camera path) under controlled conditions.
They capture interactions between components but require careful setup to ensure reproducibility.
    \item[\textbf{Application-level benchmarks}] measure end-to-end user-perceived performance (e.g.\ time from page load to interactive map) in a realistic deployment.
They are the most representative but also the noisiest due to network variability, browser behavior, and hardware differences.
\end{description}

\noindent We primarily use macro-benchmarks: we load a full style, fetch and decode tiles, and exercise the renderer with a scripted camera animation.
Our benchmark harness controls for network variability (via a caching proxy) and \ac{GPU} variability (via hardware-accelerated rendering with ANGLE/Vulkan) to achieve reproducible results at the macro level.

\begin{chaptersummary}
Three theoretical traditions map onto the three levels of the pipeline.
\textbf{Compiler optimization} (peephole rewriting, constant folding, fixpoint iteration) drives the expression-level passes that simplify each style's \ac{JSON} expressions without schema knowledge.
\textbf{Query optimization} (projection pushdown, selectivity estimation) shapes the structural passes and the pruning advisory that bridges style and data.
The \textbf{encoding techniques} surveyed last feed the data-level passes that re-encode shaved tiles with automatic strategy selection.
These cover columnar storage, the logical encodings (zigzag, delta, \ac{RLE}, dictionary) and physical techniques (VarInt, FastPFOR, \ac{FSST}) they hand off to, and space-filling curves.
We adopt the benchmarking literature's statistical methodology (bootstrap confidence intervals, warm-up handling, metric separation) rather than extending it.

The following chapter formalises the joint optimization problem (\cref{sec:problem_statement,sec:scope}) and turns these foundations into concrete passes (\cref{sec:approach,sec:building_optimizations}) and an evaluation methodology (\cref{sec:gathering_sample,sec:evaluation}).
\end{chaptersummary}
